\documentclass[10pt,a4paper]{amsart}
\usepackage[margin=19mm]{geometry}
\usepackage{amsmath,amssymb,mathtools}
\usepackage[scr=rsfs,cal=euler]{mathalfa}
\usepackage{xfrac}
\usepackage[unicode,hidelinks,bookmarks=false]{hyperref}
\newcommand{\ie}{i.e.\ }
\newcommand{\cf}{cf.\ }
\newcommand{\resp}{respectively\ }
\newcommand{\Subsection}{\subsection}
\providecommand{\nobreakdash}{\nobreak\hbox{-}}
\setlength{\emergencystretch}{2em}
\multlinegap0pt
\usepackage{mathtools}
\usepackage{amssymb}
\usepackage[shortlabels]{enumitem}
\usepackage{xcolor}
\newtheorem{lemma}{Lemma}
\newcommand{\Z}{\mathbb Z}
\title{The real quotient of the Cartwright--Steger surface}
\author{Andr\'as I. Stipsicz, Zolt\'an Szab\'o}
\date{Source: arXiv:2609.09207v1 (CC BY 4.0)}
\begin{document}
\maketitle

\noindent\emph{Source and licence.} The real quotient of the Cartwright--Steger surface. Version arXiv:2609.09207v1 (CC BY 4.0), distributed by arXiv under the Creative Commons Attribution 4.0 International licence, http://creativecommons.org/licenses/by/4.0/, as recorded in the arXiv licence metadata for this exact version and shown on the abstract page https://arxiv.org/abs/2609.09207v1. Full licence notice: \texttt{licenses/arxiv-2609.09207-CC-BY-4.0.txt}.\par\smallskip

\noindent\textit{Editorial scope.} Counted unit: the article's lemma lem:augmentation (source0.tex lines 1012--1017). The article's complete original proof of it is delivered with it (source0.tex lines 1019--1066, one \texttt{proof} environment of 48 lines). No mathematics is written, repaired or re-derived: the statement and the proof are the article's own lines, in the article's own notation, and no retained line is substituted or re-flowed.  Retained uncounted as the source-local context the proof needs: lines 937--939.  Nothing was omitted from any retained span, so the package carries no omission record.  No retained line contains a citation, so the wrapper carries no bibliography and no bibliography entry.  No retained line contains a figure environment, an image-inclusion command or drawing code, so no image is delivered and the package declares no asset dependency.  Editorial formatting only, in addition: an AMS \texttt{amsart} preamble that restates the article's own declarations and macros used by the retained lines, namely the environments \texttt{lemma} on separate counters, together with the standard packages those lines need.  The retained environments print in this document as Lemma 1 in order of appearance, whereas the article prints the same items with the numbering of its own full document.  The complete native source of the article as distributed in the arXiv e-print for this exact version is bundled unchanged as \texttt{sources/209789/source0.tex}.
\begin{equation}\label{eq:lambda-def}
 \lambda_Y(x,y)=\sum_{n\in\Z}\bigl(x\cdot t^ny\bigr)\,t^n
\end{equation}

\begin{lemma}\label{lem:augmentation}
  The natural map
  $H_2(Y;\Lambda)\otimes_{\Lambda,\varepsilon}\Z\to H_2(Y;\Z)$ is an
  isomorphism, and it carries $\lambda_Y$ to the ordinary intersection
  form $Q_Y$.
\end{lemma}

\begin{proof}
  Give $Y$ a finite CW structure and let $C_*(\widetilde Y)$ be the
  cellular chain complex of the universal cover ${\widetilde {Y}}$, a
  finite complex of free $\Lambda$-modules with $H_*(C_*(\widetilde
  Y))=H_*(Y;\Lambda)$ and $C_*(\widetilde
  Y)\otimes_{\Lambda,\varepsilon}\Z=C_*(Y)$.  Multiplication by $t-1$
  is injective on $\Lambda$ with cokernel $\Z$. As $C_*(\widetilde Y)$
  is free, with the covering map $p\colon {\widetilde {Y}}\to Y$
  inducing $p_{\#}$ on chains, we have that
\[
 0\longrightarrow C_*(\widetilde Y)\xrightarrow{\ t-1\ }C_* (\widetilde Y)
 \xrightarrow{\ p_\#\ }C_*(Y)\longrightarrow0
\]
is a short exact sequence of chain complexes.  Its homology long exact
sequence is
\[
 \cdots\to H_n(Y;\Lambda)\xrightarrow{\ t-1\ }H_n(Y;\Lambda)
 \xrightarrow{\ p_*\ }H_n(Y;\Z)\to H_{n-1}(Y;\Lambda)
 \xrightarrow{\ t-1\ }H_{n-1}(Y;\Lambda)\to\cdots .
\]
Since $H_1(Y;\Lambda)=H_1(\widetilde Y;\Z)=0$, the portion at $n=2$ reads
\[
 H_2(Y;\Lambda)\xrightarrow{\ t-1\ }H_2(Y;\Lambda)
 \xrightarrow{\ p_*\ }H_2(Y;\Z)\longrightarrow0 ,
\]
so $p_*$ is onto with kernel $(t-1)H_2(Y;\Lambda)$ and therefore induces an
isomorphism
\[
 H_2(Y;\Lambda)\otimes_{\Lambda,\varepsilon}\Z
 =H_2(Y;\Lambda)/(t-1)H_2(Y;\Lambda)
 \xrightarrow{\ \cong\ }H_2(Y;\Z),
\]
which is the map of the statement.

For the pairings, let $x,y\in H_2(Y;\Lambda)=H_2(\widetilde Y;\Z)$ be
represented by compact cycles with $x$ transverse to $t^ny$ for every $n$.
The covering $p$ is a local orientation-preserving homeomorphism, and two
points of $\widetilde Y$ have the same image if and only if they differ by a
deck translation.  Hence the transverse intersections of $p_*x$ and $p_*y$
in $Y$ correspond bijectively, and with matching signs, to the pairs
consisting of an integer $n$ and a point of $x\cap t^ny$.  Counting with
signs,
\[
 p_*x\cdot p_*y=\sum_{n\in\Z}x\cdot t^ny
 =\varepsilon\bigl(\lambda_Y(x,y)\bigr)
\]
by \eqref{eq:lambda-def}, which is the assertion about the forms.
\end{proof}

\end{document}
